
\usepackage[T1]{fontenc}
\usepackage{lmodern}

% Chinese (Simplified) support for the translated deck. xeCJK routes CJK
% codepoints to a CJK font while leaving Latin text/math on the fontspec
% faces set per-deck; this is lighter-touch than full `ctex`, which fights
% beamer/metropolis's own fontspec setup. Requires XeLaTeX (already the
% engine for these decks) and Noto CJK SC (installed system-wide).
\usepackage{xeCJK}
\setCJKmainfont{Noto Sans CJK SC}
\setCJKsansfont{Noto Sans CJK SC}
\setCJKmonofont{Noto Sans Mono CJK SC}

% Header bar matches the rafael-sharon toronto-may-2026 reference deck:
% miniframes navigation with subsection ticks, dark-teal section/subsection
% bars at the top of every content frame.
\definecolor{bg_subsec}{RGB}{43, 66, 70}
\definecolor{bg_sec}{RGB}{51, 77, 83}

\usetheme{metropolis}
\useoutertheme[subsection=true]{miniframes}
\setbeamercolor{subsection in head/foot}{fg=white, bg=bg_subsec}
\setbeamercolor{section in head/foot}{fg=white, bg=bg_sec}

\usepackage{appendixnumberbeamer}

\usepackage{amsmath, amssymb, amsthm}
\usepackage{mathtools}
\usepackage{bm}
\usepackage{graphicx}
\usepackage{booktabs}

% Code listings
\usepackage{listings}
\usepackage{xcolor}
\definecolor{jl-keyword}{HTML}{8959A8}
\definecolor{jl-string}{HTML}{718C00}
\definecolor{jl-comment}{HTML}{8E908C}
\lstdefinelanguage{Julia}{
  keywords={if,else,elseif,end,for,while,function,return,using,import,
            module,struct,abstract,type,export,begin,let,do,in,true,false,
            const,where},
  morecomment=[l]{\#},
  morestring=[b]",
  sensitive=true,
}
\lstset{
  basicstyle=\ttfamily\footnotesize,
  language=Julia,
  keywordstyle=\color{jl-keyword}\bfseries,
  stringstyle=\color{jl-string},
  commentstyle=\color{jl-comment},
  showstringspaces=false,
  breaklines=true,
  tabsize=2,
  captionpos=b,
  frame=single,
  framerule=0pt,
  backgroundcolor=\color{black!3},
  extendedchars=true,
  literate=
    {β}{{$\beta$}}1 {α}{{$\alpha$}}1 {γ}{{$\gamma$}}1 {δ}{{$\delta$}}1
    {ε}{{$\varepsilon$}}1 {ζ}{{$\zeta$}}1 {η}{{$\eta$}}1 {θ}{{$\theta$}}1
    {κ}{{$\kappa$}}1 {λ}{{$\lambda$}}1 {μ}{{$\mu$}}1 {ν}{{$\nu$}}1
    {ξ}{{$\xi$}}1 {π}{{$\pi$}}1 {ρ}{{$\rho$}}1 {σ}{{$\sigma$}}1
    {τ}{{$\tau$}}1 {φ}{{$\varphi$}}1 {χ}{{$\chi$}}1 {ψ}{{$\psi$}}1
    {ω}{{$\omega$}}1 {Λ}{{$\Lambda$}}1 {Δ}{{$\Delta$}}1 {Σ}{{$\Sigma$}}1,
}

% Common math macros (kept in sync with paper/main.tex)
\newcommand{\R}{\mathbb{R}}
\newcommand{\E}{\mathbb{E}}
\newcommand{\Prob}{\mathbb{P}}
\newcommand{\norm}[1]{\left\| #1 \right\|}
\newcommand{\Vstar}{V^{\star}}

\newcommand{\p}[1]{\left( #1 \right)}
\newcommand{\psq}[1]{\left[ #1 \right]}

% Section page: use the long-form section name (\insertsection) instead of
% the short-form (\insertsectionhead) so \section[short]{long} renders
% `long` on the section-title slide and `short` in the header. Cloned
% from metropolis's `progressbar` section page template.
\setbeamertemplate{section page}{%
  \centering%
  \begin{minipage}{22em}%
    \raggedright%
    \usebeamercolor[fg]{section title}%
    \usebeamerfont{section title}%
    \insertsection\\[-1ex]%
    \usebeamertemplate*{progress bar in section page}%
    \par%
    \ifx\insertsubsection\@empty\else%
      \usebeamercolor[fg]{subsection title}%
      \usebeamerfont{subsection title}%
      \insertsubsection%
    \fi%
  \end{minipage}%
  \par%
  \vspace{\baselineskip}%
}

% Project-specific macros
\newcommand{\Vstart}{V^{\mathrm{start}}}
\newcommand{\Vend}{V^{\mathrm{end}}}
\newcommand{\Vstarttil}{\widetilde{V}^{\mathrm{start}}}
\newcommand{\Vendtil}{\widetilde{V}^{\mathrm{end}}}
\newcommand{\Vstarthat}{\widehat{V}^{\mathrm{start}}}
\newcommand{\Vendhat}{\widehat{V}^{\mathrm{end}}}
\newcommand{\Lstart}{\Lambda^{\mathrm{start}}}
\newcommand{\Lend}{\Lambda^{\mathrm{end}}}
\newcommand{\xstart}{x^{\mathrm{start}}}
\newcommand{\xend}{x^{\mathrm{end}}}
\newcommand{\CV}{\mathcal{V}}
\newcommand{\CL}{\mathcal{L}}
